\documentclass[a4paper,12pt]{article}
\usepackage[utf8]{inputenc}
\usepackage[ngerman]{babel}
\usepackage{amsmath}
\usepackage{parskip}

\setlength{\parindent}{0pt}

\title{Modul 293 - L\"osungsblatt 13}
\author{von Lukas Meier}
\date{Unterricht vom 31.03.2026}

\begin{document}

\maketitle

\section*{L\"osungsblatt: Responsive Grundlagen}

\subsection*{Teil 1: Begriffe verstehen}
\begin{itemize}
  \item \textbf{Responsive} bedeutet, dass sich eine Webseite automatisch an verschiedene Bildschirmgr\"ossen anpasst (z.\,B. Handy, Tablet, Desktop).
  \item Die Meta-Zeile \texttt{viewport} sorgt daf\"ur, dass die Seite auf mobilen Ger\"aten in der richtigen Breite und Skalierung dargestellt wird.
  \item Fixe Pixelbreiten sind auf kleinen Ger\"aten problematisch, weil Inhalte dann zu breit werden, horizontal scrollen oder abgeschnitten wirken.
\end{itemize}

\subsection*{Teil 2: Media Queries lesen}
Gegebene Regel:
\begin{verbatim}
@media (max-width: 700px) {
  nav a {
    display: block;
  }
}
\end{verbatim}

\begin{itemize}
  \item Die Regel gilt, wenn die Viewport-Breite \textbf{700px oder kleiner} ist.
  \item Die Links im \texttt{nav} werden als Blockelemente dargestellt und stehen dadurch untereinander.
  \item Sinnvoll ist das z.\,B. f\"ur Smartphones oder schmale Tablet-Ansichten.
\end{itemize}

\subsection*{Teil 3: max-width und min-width}
\begin{itemize}
  \item \texttt{@media (max-width: 700px)}: Regel gilt bis maximal 700px (typisch f\"ur kleine Bildschirme).
  \item \texttt{@media (min-width: 700px)}: Regel gilt ab 700px (typisch f\"ur gr\"ossere Bildschirme).
\end{itemize}

Beispiels\"atze:
\begin{itemize}
  \item \texttt{max-width}: \glqq Wenn das Fenster schmal ist, wird die Navigation untereinander angezeigt.\grqq{}
  \item \texttt{min-width}: \glqq Wenn genug Platz vorhanden ist, stehen die Inhalte nebeneinander.\grqq{}
\end{itemize}

\subsection*{Teil 4: Prozentbreiten anwenden}
M\"ogliche L\"osung:
\begin{verbatim}
<!DOCTYPE html>
<html lang="de">
<head>
  <meta charset="UTF-8">
  <meta name="viewport" content="width=device-width, initial-scale=1.0">
  <title>Responsive Teil 4</title>
  <link rel="stylesheet" href="styles.css">
</head>
<body>
  <div class="wrapper">
    <section class="left">Linker Bereich</section>
    <section class="right">Rechter Bereich</section>
  </div>
</body>
</html>
\end{verbatim}

\begin{verbatim}
.wrapper {
  width: 90%;
  margin: 0 auto;
  font-size: 0;
}

.left,
.right {
  width: 48%;
  display: inline-block;
  vertical-align: top;
  font-size: 16px;
  box-sizing: border-box;
  padding: 12px;
  border: 1px solid #999;
}

.right {
  margin-left: 4%;
}
\end{verbatim}

\subsection*{Teil 5: Layout f\"ur kleine Bildschirme anpassen}
Erg\"anzung mit Media Query:
\begin{verbatim}
@media (max-width: 700px) {
  .left,
  .right {
    width: 100%;
    display: block;
    margin-left: 0;
  }

  .right {
    margin-top: 12px;
  }
}
\end{verbatim}

Ergebnis: Auf kleinen Bildschirmen werden die beiden Bereiche untereinander angezeigt.

\subsection*{Teil 6: Code erg\"anzen}
Eine passende Erg\"anzung ist:
\begin{verbatim}
.card {
  width: 45%;
  display: inline-block;
}

@media (max-width: 700px) {
  .card {
    width: 100%;
    display: block;
  }
}
\end{verbatim}

\subsection*{Teil 7: Best Practices reflektieren}
\begin{itemize}
  \item Ein Hauptbereich bekommt immer \texttt{width: 1200px}. \textbf{Eher ung\"unstig}, weil auf kleinen Ger\"aten oft horizontaler Scroll entsteht.
  \item Buttons und Texte werden auf dem Handy zus\"atzlich getestet. \textbf{Sinnvoll}, weil Bedienbarkeit und Lesbarkeit mobil gepr\"uft werden.
  \item Ein Container bekommt \texttt{width: 90\%} und \texttt{margin: 0 auto}. \textbf{Sinnvoll}, da die Breite flexibel bleibt und zentriert wird.
  \item Das Layout wird nur auf einem grossen Monitor angeschaut. \textbf{Eher ung\"unstig}, weil mobile Darstellungen unbeachtet bleiben.
\end{itemize}

\subsection*{Teil 8: Praxisaufgabe (Beispiel \glqq Mein Wochenplan\grqq{})}
\textbf{index.html}
\begin{verbatim}
<!DOCTYPE html>
<html lang="de">
<head>
  <meta charset="UTF-8">
  <meta name="viewport" content="width=device-width, initial-scale=1.0">
  <title>Mein Wochenplan</title>
  <link rel="stylesheet" href="styles.css">
</head>
<body>
  <div class="container">
    <header>
      <h1>Mein Wochenplan</h1>
    </header>

    <main class="grid">
      <section class="box">
        <h2>Schule</h2>
        <p>Montag bis Freitag: Unterricht und Projekte.</p>
      </section>

      <section class="box">
        <h2>Freizeit</h2>
        <p>Sport, Lesen und Zeit mit Freunden.</p>
      </section>
    </main>
  </div>
</body>
</html>
\end{verbatim}

\textbf{styles.css}
\begin{verbatim}
* {
  box-sizing: border-box;
}

body {
  margin: 0;
  font-family: Arial, sans-serif;
}

.container {
  width: 90%;
  margin: 20px auto;
}

.grid {
  font-size: 0;
}

.box {
  width: 48%;
  display: inline-block;
  vertical-align: top;
  font-size: 16px;
  padding: 16px;
  border: 1px solid #bbb;
  background: #f7f7f7;
}

.box + .box {
  margin-left: 4%;
}

@media (max-width: 700px) {
  .box {
    width: 100%;
    display: block;
  }

  .box + .box {
    margin-left: 0;
    margin-top: 12px;
  }
}
\end{verbatim}

\subsection*{Teil 9: Reflexion (Beispielantworten)}
\begin{itemize}
  \item Neu war f\"ur mich, wie stark \texttt{max-width} und \texttt{min-width} das Layout steuern.
  \item Gut funktioniert hat das Umschalten von nebeneinander auf untereinander mit einer Media Query.
  \item Mehr \"Ubung brauche ich bei der Wahl sinnvoller Breakpoints und beim Testen auf verschiedenen Ger\"aten.
\end{itemize}

\end{document}
